% =============================================================================
% METHODOLOGY ADDITIONS, 2026-09-24 -- DRAFT, NOT YET MERGED INTO methodology.tex
% =============================================================================
% Source of every claim: SESSION_CHANGELOG_2026-09-24.md (findings + commands),
% SFT_DATA_PROCESS.txt section 9, MCTS_SEARCH_SCORING.md "Scoring review,
% 2026-09-24". Blocks are keyed to the existing section labels; each says where
% it goes. Numbers marked [CONFIRM] depend on jobs still running (3945722/37
% rust, 3945738/44 java, 3945991 -> 3945992/93); replace them when those finish.
%
% BEFORE MERGING -- consistency fix needed in the existing text:
%   Sec. subsec:evaluator ("Fallback for trees with no valid vote") describes
%   the CURRENT rule (fallback only when V is empty), but Sec.
%   subsec:experimental_setting reports 0.9454 (java) and 0.8739 (rust), which
%   were computed with the PRE-2026-09-12 rule (override applied to every tree
%   showing the signature). Under the rule the text describes, the same trees
%   score 0.9163 and 0.7952. Either restate the numbers or describe the rule
%   that produced them (Block E gives both).
% =============================================================================


% -----------------------------------------------------------------------------
% BLOCK A -- append to \subsection{Execution Model} (subsec:execution_model)
% -----------------------------------------------------------------------------
\paragraph{Offline differential testing.} Independently of the policy model,
we run both programs of a pair on the same inputs and compare their standard
output after whitespace normalization, counting an input only when both
programs exit cleanly. Inputs come from three sources: the inputs the model
itself wrote in its test code, random inputs generated from the input
declaration of the non-Python program when it uses a declarative reader (the
declaration fixes the token sequence; several line layouts are tried because
the Python side typically reads line by line), and the official sample
inputs of the underlying programming-contest problem. For offline testing
only, the third-party libraries that competitive-programming Rust solutions
commonly depend on are compiled once and linked, so programs that the search
environment cannot run can still be executed here. This check is used to
audit labels and trajectories and to build training data
(Section~\ref{subsec:training_data}); it never takes part in producing a
prediction. On the 81 clone pairs the base model misclassifies in our Rust
training split, 60 agree on every valid input; every label-contradicting
disagreement we inspected was caused by an input that violated the problem's
constraints.

\paragraph{External-library detection.} The rule that decides whether a
non-Python program depends on unavailable libraries works on import
statements. We found that it over-matches (an import-like phrase inside a
comment, a locally declared module, or a fully qualified standard-library
path), so some Rust programs that compile without any external dependency
are reported as unrunnable. All results in this paper use the rule as
described; correcting it changes the execution environment and therefore
requires re-running every baseline.


% -----------------------------------------------------------------------------
% BLOCK B -- append to \subsection{MCTS-Based Search} (subsec:mcts_search)
% -----------------------------------------------------------------------------
\paragraph{Effective search budget.} Under the evaluation configuration
(2 rollouts, 2 candidate continuations per expansion), the requirement that
every unvisited child be visited once before any sibling is revisited means
both rollouts are spent on the two children of the root. Intermediate step
rewards therefore cannot redirect an evaluation run; they matter only at
larger budgets, such as the data-mining configurations
(Section~\ref{subsec:training_data}).

\paragraph{Step scoring.} For those larger budgets we define a label-free
step score (version 2) that fixes several properties of the earlier
consistency-based score: a step's outcome is taken from the interpreter's
exception status rather than from a substring test; a step that produces no
output, or only prints the model's own verdict, is scored neutrally rather
than positively; a passing or failing assertion is treated as evidence about
equivalence only if the executed program actually ran the non-Python side,
because otherwise it compares the Python output against a value the model
guessed; a step that runs both programs receives an additional bonus; and a
terminal clone decision whose stated reason claims tests passed, while no
step on its path produced any output, is penalized. Replaying both versions
over stored trees, and measuring how well a step's score separates steps
whose subtrees end mostly in correct answers (ground truth used only for
this measurement), the earlier score is anti-predictive on ground-truth
clone questions (AUC 0.35--0.48), and version 2 reverses this on mining
trees (AUC 0.57--0.61). Both stay close to chance overall, so we treat step
scoring as a weak lever and do not use version 2 for evaluation.


% -----------------------------------------------------------------------------
% BLOCK C -- insert into \subsection{The Evaluator} (subsec:evaluator),
%            after "Vote aggregation"
% -----------------------------------------------------------------------------
\paragraph{Test-conditioned vote correction.} The prompt instructs the model
to skip testing and emit a fixed non-clone marker when it judges a pair
non-equivalent, and to test otherwise (Section~\ref{subsec:prompt_design}).
A non-clone leaf whose trajectory nevertheless ran a test is therefore a
trajectory that leaned towards equivalence and reversed its decision after a
test -- in practice, usually a test comparing the Python output against an
expected value the model guessed for the other program. We therefore let
such a leaf vote \texttt{clone}:
\[
    v_i' =
    \begin{cases}
        \text{clone}, & \text{if } v_i = \text{non-clone} \text{ and trajectory } i \text{ executed a test}, \\
        v_i, & \text{otherwise},
    \end{cases}
\]
before applying clone-on-disagreement to $V' = \{v_1', \ldots, v_k'\}$.
Like clone-on-disagreement, this correction depends on the model's
behaviour rather than on the task: it is justified only when non-clone
decisions reached after testing are mostly wrong, which holds for the base
model (91--96\% of such leaves are ground-truth clones on the training split)
but not for our fine-tuned models, which also test genuinely non-equivalent
pairs. We therefore decide per model whether to apply it, on that model's own
predictions for the training split, and report only held-out numbers.


% -----------------------------------------------------------------------------
% BLOCK D -- new \subsection{Training Data Construction}, before
%            \subsection{Experimental Setting}
% -----------------------------------------------------------------------------
\subsection{Training Data Construction}
\label{subsec:training_data}

\paragraph{Splits.} Fine-tuning data drawn from the evaluation benchmark is
restricted to a training split, and models trained on it are evaluated only
on the corresponding held-out split. Splitting by the Python side's problem
alone is insufficient: the benchmark reuses each non-Python program in its
clone pair and in several non-clone pairs, and under a Python-only split
every held-out Rust program (and 19\% of held-out Java programs) also appears
in the training split. We therefore assign each non-Python program to the
problem of the clone pair it appears in, draw 20\% of problems as held out,
and keep a pair on a side only if both of its programs' problems belong to
that side, dropping pairs that straddle the split. For Rust this yields 830
training and 110 held-out pairs (60 clone, 50 non-clone), with no exact and
no near-duplicate ($\geq$0.95 token Jaccard) programs shared between them.

\paragraph{Trajectory filtering.} Trajectories mined from the model's own
search are kept for supervised fine-tuning only if they reach the correct
label and their evidence is sound. A trajectory is discarded if it attempted
execution and never obtained a successful run, if its answer claims that
tests passed while no step produced output, if it concluded equivalence from
an input on which the two programs actually differ, or if its only
``test'' printed a verdict. Trajectories that are correct but unverified
(for example, whose only code step produced no output) are kept at lower
priority, and near-duplicate trajectories are capped per question. We then
reviewed every remaining trajectory manually and removed those that describe
the wrong problem, test a re-implementation rather than the given program,
state false facts about either program, or reach the right label with a
wrong description.

\paragraph{Teacher and repaired trajectories.} For misclassified clone pairs
that the model's own search never solves, we add trajectories whose analysis
steps were written by hand after reading both programs, in the same step
format the policy produces. Their test steps run both programs, or, when the
execution environment cannot run the non-Python side, the Python side
against that side's verified outputs, on inputs where the two programs were
verified to agree. The observation shown in every such trajectory is the
real output of the execution environment, never written by hand, and a
trajectory is kept only if every compared case matches. Where a model
trajectory's reasoning was correct but its test step produced nothing, we
keep its reasoning and conclusion verbatim and replace only the test step.
Because these trajectories are written by a different model, fine-tuning on
them is a form of distillation, and we report it separately from
self-generated data. Non-clone examples are drawn from the model's own
verified trajectories so that both labels are equally represented.

\paragraph{Preference data.} From the same trees we derive step-level
preference pairs (sibling steps whose subtrees differ in mean trajectory
reward by at least 0.5) and trajectory-level pairs (best- vs.
worst-rewarded trajectory of a question, a teacher trajectory vs. the model's
most-visited wrong one, and a repaired trajectory vs. its original). The
trajectory reward is $+1$ for a correct decision backed by real execution
output, $+0.5$ for a correct but unverified one, $-0.5$ for a correct
decision whose trajectory has one of the defects listed above, and $-1$ for
a wrong decision, so that correct-for-the-wrong-reason trajectories are not
reinforced.


% -----------------------------------------------------------------------------
% BLOCK E -- \subsection{Experimental Setting} (subsec:experimental_setting):
%            replace the baseline numbers / add to "Metrics"
% -----------------------------------------------------------------------------
\paragraph{Aggregation version.} All $F_1$ scores are computed from the
stored search trees with a single, fixed evaluator version, and every model
is re-scored with that version. [Choose one and state it:]
% OPTION 1 -- current rule (fallback only when no leaf votes):
%   base Qwen3-4B full CLCCD: Python--Java 0.9163, Python--Rust 0.7952;
%   codenet-all-v1: Java 0.9425, Rust 0.8498.
% OPTION 2 -- pre-2026-09-12 rule (signature overrides every tree):
%   base: Java 0.9454 (0.9443 recomputed), Rust 0.8739;
%   codenet-all-v1: Java 0.8594, Rust 0.6354 (the old rule collapses
%   fine-tuned models' precision, 0.51 on Rust).
% Held-out numbers (training/held-out split, Block D), base Qwen3-4B:
%   Rust rl_heldout: current 0.8039, with test-conditioned correction 0.8909
%   (independent second run 0.7921 -> 0.8850);
%   Rust rl_heldout_strict: 0.8000 -> 0.8785 (second run 0.7174 -> 0.8762);
%   Java rl_heldout: 0.9080 -> 0.9634.
%   [CONFIRM with 3945722/3945738 _result_NEWRULE_tested_agg]


% -----------------------------------------------------------------------------
% BLOCK F -- append to \subsection{Limitations} (subsec:limitations)
% -----------------------------------------------------------------------------
Two evaluator components, clone-on-disagreement and the test-conditioned vote
correction, exploit asymmetries in a particular model's behaviour rather than
properties of the task; we decide whether to apply each on training-split
predictions of the same model and report held-out results, but their gains
should not be expected to transfer to other models without re-validation.
Our strictest held-out Rust split contains only 110 pairs, so differences of
a few hundredths of $F_1$ on it are within run-to-run variation (about $\pm$0.01
between independent runs of the same model, larger on that subset).
Fine-tuning data that includes teacher-written trajectories measures the
effect of distillation, not of self-improvement alone.
